\documentclass[12pt,conference,onecolumn]{IEEEtran}

\usepackage[hidelinks]{hyperref}

\title{Comparative fluid-dynamic analysis of {B}eloniform fishes and the development of a biomimetic robotic model}
\author{%
\IEEEauthorblockN{Mark Reznik}\IEEEauthorblockA{Science \& Engineering\\Manalapan High School\\Englishtown, NJ\\\href{mailto:427mreznik@frhsd.com}{427mreznik@frhsd.com}}\and
\IEEEauthorblockN{Soham Sankritya}\IEEEauthorblockA{Science \& Engineering\\Manalapan High School\\Englishtown, NJ\\\href{mailto:427ssanrkitya@frhsd.com}{427ssankritya@frhsd.com}}\and
\IEEEauthorblockN{Lorenzo Brunie}\IEEEauthorblockA{Science \& Engineering\\Manalapan High School\\Englishtown, NJ\\\href{mailto:427lbrunie@frhsd.com}{427lbrunie@frhsd.com}}
}

\date{January 28, 2027}

\newcommand{\keywords}{biomechanics}

\usepackage{hyperref}
\makeatletter
\AtBeginDocument{
\hypersetup{%
pdftitle={\@title},
pdfauthor={Mark Reznik, Soham Sankritya, and Lorenzo Brunie},
pdfkeywords={\keywords}}}
\makeatother

\begin{document}
\maketitle 

\begin{abstract}
Fish have existed for hundreds of millions of years and have developed into a variety of morphologies and swimming modes that can help to better our understanding of fluid dynamics and biomechanics. Two surface dwelling families within the order Beloniformes, the needlefishes (Belonidae) and flying fishes (Exocoetidae), order Beloniformes, share similar habitats but exhibit very different morphologies and swimming modes. In this project, the fluid dynamics of both families will be compared with a focus on how morphological traits affect the swimming mode of each family. Needlefish are reported to utilize anguilliform (eel-like) swimming along with the use of fins. For this project, we will use the needlefish as a model for our robotic system. The robot will not aim to mimic the external morphology of the fish but rather specific morphological traits that contribute to the needlefish’ swimming mode. The robot will consist of a controller (e.g. Arduino), motor or servo and a tail or linkage. The body of the robot will be designed in a CAD program and printed in 3D. The tail will be made of flexible material and sealed to prevent water intrusion. The robot will be tested in a water tank measuring displacement, swimming speed and thrust. The goal of this project is to demonstrate the ability to create a functional biomimetic robot with a trait observed in nature.
\end{abstract}

\begin{IEEEkeywords}
\keywords
\end{IEEEkeywords}

\end{document}
